\providecommand{\main}{../main}
\documentclass[\main/main.tex]{subfiles}
\graphicspath{{../images/}}
\begin{document}
%------------------------------------------------%
\section{
    共形不変性
}
%------------------------------------------------%
この章では共形変換を具体的に定義して、共形変換に対して良い変換性を示すプライマリー場を導入する。
以降の議論では常に平坦な時空の計量はEuclid計量$δ_{μν}$とするが、これを全てMinkowski計量$η_{μν}$に置き換えてもほとんど同じ議論が可能である。
\subsection{
    Euclid群
}
共形変換群を定義する前に、その部分群であるEuclid群について見ておく。これは並進と回転から構成される群である。

$d$次元空間における線素は計量テンソル$g_{μν}$を用いて、
\begin{align}
    \d{s}² = g_{μν}(x)\d{x^μ}\d{x^{ν}}
\end{align}
と表せる。
座標変換$x → x'$に対し、$\d{{s'}^2}=\d{s^2}$を課すことで、計量の変換$g_{μν}(x) → g'_{μν}(x')$が得られる。
微小変換$x^μ → x'^μ = x^μ - ε^μ (x)$に対し、この変換は
\begin{align}
    g'_{μν}(x')
    = \∂{x^ρ}{x'^μ}\∂{x_σ}{x'^ν}g_{ρσ}(x)
    = g_{μν}(x) + ∂_με_ν(x) + ∂_νε_μ(x)
    \label{eq: infinitesimal metric transformation}
\end{align}
で与えられる。

以降はEuclid空間を考え、$g_{μν}(x) = δ_{μν}$と固定する。
計量を不変に保つような座標変換は群をなし、(特殊)Euclid群と呼ばれる。
微小変換がEuclid群に属する条件は、
(\ref{eq: infinitesimal metric transformation})
からKilling方程式
\begin{align}\tcboxmath{
    δg_{μν} = ∂_μ ε_ν +  ∂_ν ε_μ = 0
    \label{killing eq}
}\end{align}
を満たすことである。
この方程式を解いてみよう。
(\ref{killing eq})に$∂_ρ$を作用させると、
\begin{align}
    ∂_μ∂_ρε_ν + ∂_ν∂_ρε_μ
    = -2 ∂_μ ∂_ν ε_ρ = 0
\end{align}
となる。
ただし1つ目の等号で(\ref{killing eq})を利用した。このことから$ ε_μ$は高々$x$の一次式である。すなわち、
\begin{align}
     ε_μ = a_μ + b_{μν}x^{ν}
\end{align}
と表せる。
これを(\ref{killing eq})に代入することで、$b_{μν}$が反対称であることがわかる。
ここで、$x^μ-ε^μ(x)=(1-ε^ν(x)∂_ν)x^μ$と書くことで、$ε^ν(x)∂_ν$を座標変換の生成子とみなせる。
並進と回転を生成する微分演算子は
\begin{align}
    &
    𝓅_μ ≔ ∂_μ,
    \\ &
    𝓂_{μν} ≔ -x_μ∂_ν+x_ν∂_μ
\end{align}
と定義される。
これらの間の交換関係を計算すると、
\begin{align}
    &
    [𝓅_μ, 𝓅_ν] = 0,
    \\ &
    [𝓅_μ, 𝓂_{νρ}]
    = -δ_{μν}𝓅_ρ+δ_{μρ}𝓅_ν,
    \label{comm. relation p-m}
    \\ &
    [𝓂_{μν}, 𝓂_{ρσ}]
    = -δ_{νρ}𝓂_{μσ}+δ_{μρ} 𝓂_{νσ}
    +δ_{νσ}𝓂_{μρ}-δ_{μσ}𝓂_{νρ}
    \label{comm. relation m-m}
\end{align}
となる。
単純計算なので過程は載せないが、(\ref{comm. relation p-m}),(\ref{comm. relation m-m})において基本的な関係式は
\begin{align}
    &
    [𝓅_1,𝓂_{23}] = [𝓂_{12},𝓂_{34}] = 0,
    \\ &
    [𝓅_1,𝓂_{12}] = -𝓅_2,
    \\ &
    [𝓂_{12},𝓂_{23}] = -𝓂_{13}
\end{align}
であり、これを$𝓂_{μν}$の添字の反対称性に注意して拡張すれば(\ref{comm. relation p-m}),(\ref{comm. relation m-m})が得られる。
% \footnote{
% \begin{align*}
%     [𝓅_μ,𝓂_{νρ}]
%     &
%     = -∂_μ(x_ν∂_ρ-x_ρ∂_ν)+(x_ν∂_ρ-x_ρ∂_ν)∂_μ
%     \\ &
%     = -δ_{μν}𝓅_ρ+δ_{μρ}𝓅_ν
%     \\
%     [𝓂_{μν},𝓂_{ρσ}]
%     &
%     = (x_μ∂_ν - x_ν∂_μ)(x_ρ∂_σ - x_σ∂_ρ)
%     - (x_σ∂_ρ - x_ρ∂_σ)(x_ν∂_μ - x_μ∂_ν)
%     \\ &
%     = [δ_{νρ}x_μ∂_σ - (μ↔ν) - (σ↔ρ) + (μρ↔νσ)]
%     \\ &\␣\␣
%     - [δ_{νρ}x_σ∂_μ - (μ↔ν) - (σ↔ρ) + (μρ↔νσ)]
%     \\ &
%     = -δ_{νρ}𝓂_{μσ} - (μ↔ν) - (σ↔ρ) + (μρ↔νσ)
% \end{align*}
% }
Euclid群の一般の元は有限のパラメーター$a^μ,b^{μν}$と指数関数を用いて
\begin{align}
    \exp(- a^μ𝓅_μ - \f{1}{2}b^{μν}𝓂_{μν})
\end{align}
と表される。

\subsection{
    共形変換群
}
ここで条件を緩めて、計量が
\begin{align}
    g'_{μν} = Ω(x)^2 g_{μν}
\end{align}
と変換する座標変換$x → x'$を考える。
これを共形変換とよぶ。
ここで、
\begin{align}
    \∂{x^μ}{x'^ν} = Ω(x){R_ν}^μ(x),\␣
    {R_ν}^μ ∈ \rm{SO}(d)
\end{align}
と表すことができる。
したがって共形変換は局所的な角度を保つ変換である。
\begin{figure}[H]
    \centering
    \includegraphics[width=0.5\hsize]{../images/conformal transformation.pdf}
    \caption{共形変換}
\end{figure}

微小な共形変換$x^μ → {x'}^μ = x^μ -  ε^μ $に対し、
\begin{align}
     ∂_μ ε_ν +  ∂_ν ε_μ = λ(x)δ_{μν}
    \label{def: lambda}
\end{align}
と書ける。ただし$λ(x)$は任意の関数である。
両辺のトレースをとると、
\begin{align}
    λ(x) = \f{2}{d} ∂_μ ε^μ
    \label{eq: lambda-ε relation}
\end{align}
となる。これを元の式に代入して以下の共形Killing方程式を得る。
\begin{align}\tcboxmath{
    ∂_μ ε_ν+ ∂_ν ε_μ = \f{2}{d}δ_{μν} ∂_ρ ε^ρ
    \label{conformal killing eq}
}\end{align}
以下でこれを解こう。過程が気にならない方は(\ref{eq: solution of conformal killing equation})に飛んでしまって構わない。
\eqref{def: lambda}の両辺を$ ∂_ρ$で微分し、\eqref{def: lambda}を用いると、
\begin{align}
    & \nonumber
     ∂_ρ ∂_μ ε_ν +  ∂_ρ ∂_ν ε_μ
    \\ &
    = -2 ∂_μ ∂_ν ε_ρ
    +  ∂_μ λ(x) δ_{νρ}
    +  ∂_ν λ(x) δ_{μρ}
    =  ∂_ρ λ(x) δ_{μν}
    \label{eq: conformal ε condition 2}
\end{align}
となる。さらに$μ,ν$を縮約して、
\begin{align}
    -2 ∂_μ ∂^μ ε_ρ
    = (d-2) ∂_ρ λ(x)
\end{align}
を得る。この式を、\eqref{eq: conformal ε condition 2}の両辺を$ ∂^ρ$で微分した式に代入すると、
\begin{align}
     ∂^ρ∂_ρ ∂_μ ε_ν
    +  ∂^ρ∂_ρ ∂_ν ε_μ
    = (2-d) ∂_μ ∂_ν λ(x)
    =  ∂^ρ∂_ρ λ(x)  δ_{μν}.
    \label{eq: conformal condition 1}
\end{align}
また$μ,ν$を縮約して、
\begin{align}
    (1-d) ∂_μ ∂^μ λ(x) = 0.
    \label{eq: conformal condition 2}
\end{align}
$d>2$の場合、
\footnote{
    $d=2$の場合、(\ref{eq: conformal condition 1})の条件が緩むことで共形対称性は無限個の生成子を持つ。
    また$d=1$の場合は(\ref{eq: conformal condition 1}),(\ref{eq: conformal condition 2})のどちらも消える。
    これは1次元で角度が定義されないことに対応している。
}
(\ref{eq: conformal condition 1})、
(\ref{eq: conformal condition 2})から
    $ ∂_μ ∂_ν λ(x) = 0$
を得る。したがって$λ(x)$は高々1次である。
また(\ref{eq: lambda-ε relation})から$ ε_μ(x)$は高々2次であり、
\begin{align}
    ε_μ
    = a_μ + b_{μν}x^{ν} + c_{μνρ}x^νx^ρ
\end{align}
とおける。ここで$c_{μνρ}=c_{μρν}$とする。
これを(\ref{conformal killing eq})
に代入すると、1次の項に関して、
\begin{align}
    b_{μν} + b_{νμ}
    = \f{2}{d}{b^ρ}_ρ δ_{μν}
\end{align}
が得られる。ここから、$b_{μν}$の対称成分が、$δ_{μν}$に比例することが分かる。
したがって、
$b_{μν}$の対称成分が生成する変換は
$x^μ → x'^μ = (1-λ)x^μ$と書ける。
これはスケール変換になっている。

また(\ref{def: lambda})から、
2次の項に関して、
\begin{align}
    2c_{νρμ}+2c_{μρν}
    =\f{2}{d}δ_{μν}{c^σ}_{σρ}
\end{align}
となる。ここから
\begin{align}
    c_{μνρ}
    = - c_μ δ_{νρ}
    + c_ν δ_{ρμ}
    + c_ρ δ_{μν},
    \␣
    c_μ = \f{1}{d}{c^σ}_{σμ}
\end{align}
となる。
$c_{μνρ}$が生成する座標変換は、
$
    x^μ → x'^μ
    = x^μ - 2(x^{ν} c_ν)x^μ + c^μ x^2
$
であり、以上により(\ref{conformal killing eq})の解は
\begin{align}\tcboxmath{
    ε^μ = a^μ + b^{μν}x_ν + λx^μ + c^ν(2x_νx^μ-x^2δ^μ_ν)
    \label{eq: solution of conformal killing equation}
}\end{align}
と書ける。ただし$a^μ,b^{μν},λ,c^μ$は微小なパラメーターであり、$b^{νμ}=-b^{μν}$である。
ここから共形変換を生成する微分演算子を
\begin{align}
    &
    𝓅_μ ≔ ∂_μ,
    \\ &
    𝓂_{μν} ≔ -x_μ∂_ν+x_ν∂_μ,
    \\ &
    𝒹 ≔ x⋅∂,
    \\ &
    𝓀_μ ≔ 2x_μ(x⋅∂) - x² ∂_μ
\end{align}
と定義できる。$𝓅_μ,𝓂_{μν}$はEuclid群の生成子と同一であり、
新たにスケール変換の生成子$𝒹$と特殊共形変換の生成子$𝓀_μ$が追加されている。

$𝓅_μ,𝓂_{μν},𝒹$が生成する有限の共形変換はそれぞれ並進、回転、スケール変換であり、分かりやすい。
$𝓀_μ$が生成する有限の変換は特殊共形変換であり、
\begin{align}
    x^μ → {x'}^μ = \f{x^μ-x²c^μ}{1-2c⋅x+c²x²}
    \label{special conformal transformation}
\end{align}
という形になる。この変換は反転変換、並進変換、反転変換の組み合わせで構成される。
\begin{align}
    x^μ → \f{x^μ}{x²}
    → \f{x^μ}{x²} - c^μ
    → \f{x^μ/x² - c^μ}{(x/x²-c)²}
    = \f{x^μ-x²c^μ}{1-2c⋅x+c²x²}
\end{align}
特殊共形変換がパラメーター$c^μ$について加法性を満たすことは、反転変換を2回繰り返すと元に戻ることから、並進変換の加法性より明らかである。
また$c^μ$が微小だとして展開すると、
\begin{align}
    {x'}^μ
    = x^μ + 2(c⋅x)x^μ - x²c^μ
    = (1 + c⋅𝓀)x^μ
\end{align}
となるから、(\ref{special conformal transformation})は確かに$𝓀_μ$で生成される変換になっている。
特殊共形変換が反転変換、並進変換、反転変換の合成で構成されることを微小変換に適用すると、生成子の間の関係
\begin{align}
    𝓀_μ = -I𝓅_μI
\end{align}
が分かる。ただし$I$は反転を表す。

$𝓅_μ, 𝓂_{μν}, 𝒹, 𝓀_μ$の満たす交換関係を計算すると、以下のようになる。
% \footnote{
% \begin{align*}
%     [𝒹,𝓅_μ] &
%     = [x⋅∂, ∂_μ]
%     = - ∂_μ
%     = -𝓅_μ
%     \\
%     [𝒹,𝓂_{μν}] &
%     = [x·∂,-x_μ∂_ν]-(μ ↔ ν)
%     = 0
%     \\
%     [𝒹,𝓀_μ] &
%     = [x⋅∂, 2x_μ(x⋅∂)]
%     - [x·∂, x^2∂_μ]
%     = (4x_μ(x·∂)- 2x_μ(x·∂))
%     - (2x^2 ∂_μ- x^2 ∂_μ)
%     \\ &
%     = 2x_μ(x·∂)-x^2∂_μ
%     = 𝓀_μ
%     \\
%     [𝓀_μ, 𝓀_ν]
%     & = [
%         2x_μ x · ∂ - x^2  ∂_μ
%         ,
%         2x_ν x · ∂ - x^2  ∂_ν
%     ]
%     \\ &
%     = 8x_μ x_ν (x · ∂)
%     - 4x^2 x_μ  ∂_ν
%     - 2δ_{μν} (x^2 x · ∂)
%     - 2x^2 x_ν  ∂_μ
%     + 2x^2 x_μ  ∂_ν - (μ ↔ ν)
%     \\ &
%     = 0
%     \\
%     [𝓀_μ, 𝓅_ν]
%     &
%     = [2x_μ(x·∂)-x^2 ∂_μ, ∂_ν]
%     = -2δ_{μν}(x · ∂)
%     - 2x_μ  ∂_ν
%     + 2x_ν  ∂_μ
%     \\ &
%     = 2𝓂_{μν} -2δ_{μν}𝒹
%     \\
%     [𝓀_μ, 𝓂_{νρ}]
%     &
%     = -[
%         2x_μ(x·∂)-x^2 ∂_μ
%         ,
%         x_ν ∂_ρ - x_ρ ∂_ν
%     ]
%     \\  &
%     = -2x_μ x_ν  ∂_ρ
%     + x^2 δ_{μν} ∂_ρ
%     + 2 δ_{μρ} x_ν (x·∂)
%     + 2 x_μ x_ν  ∂_ρ
%     - 2 x_ν x_ρ  ∂_μ
%     - (ν ↔ ρ)
%     \\ &
%     = δ_{μν}(-x^2 ∂_ρ+2x_ρ(x·∂)) - (ν ↔ ρ)
%     \\ &
%     = -δ_{μν}𝓀_ρ + δ_{μρ}𝓀_ν
% \end{align*}
% }
\begin{align}
    &
    [𝓅_μ, 𝓅_ν] = [𝒹,𝓂_{μν}] = [𝓀_μ,𝓀_ν] = 0,
    \\ &
    [𝓂_{μν},𝓂_{ρσ}]
    = -δ_{νρ}𝓂_{μσ}+δ_{μρ} 𝓂_{νσ}+δ_{νσ}𝓂_{μρ}-δ_{μσ}𝓂_{νρ},
    \\ &
    [𝒹,𝓅_μ] = -𝓅_μ,
    \\ &
    [𝓅_μ,𝓂_{νρ}] = -δ_{μν}𝓅_ρ + δ_{μρ}𝓅_ν,
    \\ &
    [𝒹,𝓀_μ] = 𝓀_μ,
    \\ &
    [𝓀_μ,𝓂_{νρ}] = -δ_{μν}𝓀_ρ + δ_{μρ}𝓀_ν,
    \\ &
    [𝓀_μ,𝓅_ν] = 2𝓂_{μν} - 2δ_{μν}𝒹
\end{align}
このように定義されるLie代数を\textbf{($d$次元)共形代数}と呼ぶ。
$d=2$の場合、共形変換の生成子は無限次元に拡大するが、$𝓅_μ,𝓂_{μν},𝒹,𝓀_μ$はその中で部分代数(大域的共形代数)をなす。
対称性を大域的共形代数に限ることにすれば、本記事の以降の議論は$d=2$の場合も有効である。

詳細な計算過程は省くが、いくつか補足をしておく。
まず$[𝒹,𝓂_{μν}] = 0$はスケール変換と回転が可換であることを意味しており、直観的には明らかである。
% 次に$[𝒹,𝓅_μ]=-𝓅_μ$は
% \begin{align}
%     [x⋅∂,∂_μ] = -∂_μ(x^ν∂_ν) = -∂_μ
% \end{align}
% から分かる。
$𝓀_μ$についての交換関係は
\begin{align}
    𝓀_μ = -I𝓅_μI,\␣ 𝓂_{μν} = I𝓂_{μν}I, \␣ 𝒹 = -I𝒹I
\end{align}
を用いると$𝓅_μ$についての交換関係に帰着できる。
例えば、
\begin{align}
    [𝓀_μ, 𝓀_ν]
    = [I𝓅_μI,I𝓅_νI]
    = I[𝓅_μ,𝓅_ν]I = 0
\end{align}
である。$[𝓀_μ, 𝓅_ν]$についてはこのトリックが使えないので真面目な計算を載せておく。
\begin{align}
    [𝓀_μ, 𝓅_ν]
    &
    = [2x_μ(x·∂)-x^2 ∂_μ, ∂_ν]
    \∅ &
    = -2δ_{μν}(x · ∂)
    - 2x_μ  ∂_ν
    + 2x_ν  ∂_μ
    \∅ &
    = 2𝓂_{μν} -2δ_{μν}𝒹
\end{align}

\subsection{
    積の順序についての注意
}
微分演算子の積の順序について注意をしておく。
例として
\begin{align}
    \e^{-b⋅𝓂}x^μ ≕ {R^μ}_ν(b)x^ν,
    \␣
    \Q(b⋅𝓂 = \f{1}{2}b^{μν}𝓂_{μν})
\end{align}
とおくと、
\begin{align}
    \e^{-b'⋅𝓂}\e^{-b⋅𝓂}x^μ
    &
    = {R^μ}_ν(b)\e^{-b'⋅𝓂}x^ν
    \∅ &
    = {R^μ}_ν(b){R^ν}_ρ(b')x^ρ
\end{align}
となる。
すなわち、微分演算子の順序と座標変換の順序は逆である。
このことは回転変換以外についても成り立つ。
一般の微分演算子$𝒬,𝒬'$について、座標変換の合成
\begin{align}
    % x' = \e^{𝒬(x)}x,\␣
    x'' = \e^{𝒬'(x')}x' = \e^{𝒬'(x')}\e^{𝒬(x)}x
\end{align}
を考える。ここで$𝒬,𝒬'$が一般に座標$x$に依存することに注意する。$𝒬'$の引数を$x$に直すと、
\begin{align}
    𝒬'(x') = \e^{𝒬(x)}𝒬'(x)\e^{-𝒬(x)}
\end{align}
となるので
\begin{align}
    x'' = \e^{𝒬(x)}\e^{𝒬'(x)}\e^{-𝒬(x)}x'
    = \e^{𝒬(x)}\e^{𝒬'(x)}x
\end{align}
となる。
したがって座標変換の合成を微分演算子で表現する場合は積を逆順にしなければならない。

これはややこしいので、座標変換の合成と同じ順序に従う共形変換の生成子$P_μ,M_{μν},D,K_μ$を定義する。
これらの生成子の積は微分演算子の逆順の積に対応するとする。
例えば、
\begin{align}
    P_μK_ν ↔ 𝓀_ν𝓅_μ
\end{align}
とする。したがって符号を反転した以下の交換関係が成り立つ。
\begin{align}
    &
    [P_μ, P_ν] = [D,M_{μν}] = [K_μ,K_ν] = 0,
    \\ &
    [M_{μν},M_{ρσ}]
    = δ_{νρ}M_{μσ}-δ_{μρ} M_{νσ}-δ_{νσ}M_{μρ}+δ_{μσ}M_{νρ},
    \\ &
    [D,P_μ] = P_μ,
    \\ &
    [P_μ,M_{νρ}] = δ_{μν}P_ρ - δ_{μρ}P_ν,
    \\ &
    [D,K_μ] = -K_μ,
    \\ &
    [K_μ,M_{νρ}] = δ_{μν}K_ρ - δ_{μρ}K_ν,
    \\ &
    [K_μ,P_ν] = 2δ_{μν}D - 2M_{μν} 
\end{align}
以降の議論でも共形変換の生成子は頻繁に用いるが、筆記体は微分演算子を意味し、イタリックの生成子とは積の順序が逆なので、注意してほしい。
\subsection{
    プライマリー場
}
座標変換$x → x'$によって計量が
\begin{align}
    g'_{μν} = Ω(x)^2 g_{μν}
\end{align}
となったとき、場$𝒪(x)$の変換性が
\begin{align}
    𝒪'(x')
    = Ω(x)^Δ 𝒪(x)
    \label{definition of primary fields}
\end{align}
で与えられるならば、$𝒪(x)$を\textbf{共形次元}$Δ$の\textbf{プライマリー場}と呼ぶ。
空間の添字をもつ場の場合、
\begin{align}
    𝒪_{μ₁⋯μₗ}(x)\d{x^{μ₁}}⋯\d{x^{μₗ}}
\end{align}
を共形次元$Δ$のプライマリー場とみなせる。
すなわち場の変換性は
\begin{align}
    𝒪'_{μ₁⋯μₗ}(x')
    &
    = Ω(x)^{Δ-l}\∂{x^{ν₁}}{{x'}^{μ₁}}⋯\∂{x^{νₗ}}{{x'}^{μₗ}}𝒪_{ν₁⋯νₗ}(x)
    \∅ &
    = Ω(x)^Δ
        {R_{μ_1}}^{ν_1}(x) ⋯ {R_{μ_l}}^{ν_l}(x)
        𝒪_{ν₁⋯νₗ}(x)
    \label{definition of primary tensor fields}
\end{align}
で与えられる。

微小な共形変換$x^μ → x'^μ = x^μ - ε^μ$に対するプライマリー場の変換性を求めよう。
ただし
\begin{align}
    ε^μ = a^μ + b^{μν}x_ν + λx^μ + 2(c⋅x)x^μ - c^μx²
\end{align}
であり、$b_{μν}$は反対称とする。
過程が気にならない方は(\ref{conformal transformation on primary fields})に飛んでしまって構わない。
(\ref{conformal killing eq})から
\begin{align}
    Ω(x)^2 δ_{μν}
    = \Q(1 +\f{2}{d}∂⋅ε)δ_{μν}
    = \Q(1+2(λ+2c⋅x))δ_{μν}
\end{align}
となるので、
\begin{align}
    Ω(x)^Δ = 1 + λΔ + 2(c⋅x)Δ
\end{align}
が分かる。
また
\begin{align}
    \∂{x_ν}{{x'}^μ}
    &
    = δ_{μν} + ∂_με_ν
    \∅ &
    = (1+λ+2c⋅x)δ_{μν} + b_{νμ} + 2(c_μx_ν-c_νx_μ)
    \∅ &
    = Ω(x){R^ν}_μ(x)
\end{align}
より
\begin{align}
    R_{μν}(x)
    &
    = δ_{μν} + b_{νμ} + 2(c_μx_ν-c_νx_μ)
    \∅ &
    = δ_{μν} + (b_{ρσ}-2c_σx_ρ)(δ^ρ_νδ^σ_μ-δ^ρ_μδ^σ_ν)
\end{align}
となるので、
\begin{align}
    R_{μ_1ν_1}(x) ⋯ R_{μ_lν_l}(x)
    = δ_{μ_1ν_1}⋯δ_{μ_lν_l}
    + (b_{ρσ}-2c_σx_ρ)
        ∑_i(δ^ρ_{ν_i}δ^σ_{μ_i}-δ^ρ_{μ_i}δ^σ_{ν_i})
            ∏_{j ≠ i}δ_{μ_jν_j}
\end{align}
と書ける。
また
\begin{align}
    &
    𝒪'_{μ₁⋯μₗ}(x)-𝒪'_{μ₁⋯μₗ}(x')
    \∅ &
    = ε⋅∂𝒪_{μ₁⋯μₗ}(x)
    = \Q(a^μ𝓅_μ+\f{1}{2}b^{μν}𝓂_{μν}+λ𝒹+c^μ𝓀_μ)
        𝒪_{μ_1⋯μ_l}(x)
\end{align}
である。
以上の式を(\ref{definition of primary tensor fields})に代入することにより、
\begin{align}
    &
    𝒪'_{μ₁⋯μₗ}(x)
    = \Q(
        a^μ𝒫_μ
        +\f{1}{2}b^{μν}ℳ_{μν}
        +λ𝒟
        +c^μ𝒦_μ
    )𝒪_{μ₁⋯μₗ}(x)
    \label{conformal transformation on primary fields}
\end{align}
と書ける。ここで、
\begin{align}
    &
    𝒫_μ ≔ ∂_μ
    \label{action of P on primary fields}
    \\ &
    ℳ_{μν} ≔ -x_μ∂_ν+x_ν∂_μ+𝒮_{μν}
    \label{action of M on primary fields}
    \\ &
    𝒟 ≔ x⋅∂+Δ
    \label{action of D on primary fields}
    \\ &
    𝒦_μ ≔ 2x_μ(x⋅∂) - x^2∂_μ + 2x_μΔ - 2x^ν𝒮_{μν}
    \label{action of K on primary fields}
\end{align}
と定義した。また
\begin{align}
    (𝒮_{μν})_{μ₁⋯μₗν₁⋯νₗ}𝒪^{ν₁⋯νₗ}
    ≔ ∑_{j=1}^l (δ^μ_{νⱼ}δ^ν_{μⱼ}-δ^μ_{μⱼ}δ^ν_{νⱼ})
        𝒪_{μ₁⋯νⱼ⋯μₗ}
\end{align}
である。 
共形変換の生成子$P_μ,M_{μν},D,K_μ$はプライマリー場に対し、
\begin{align}
    &
    P_μ𝒪(x) = 𝒫_μ𝒪(x),
    \\ &
    M_{μν}𝒪(x) = ℳ_{μν}𝒪(x)
    \\ &
    D𝒪(x) = 𝒟𝒪(x),
    \\ &
    K_μ𝒪(x) = 𝒦_μ𝒪(x)
\end{align}
を満たす。
(\ref{action of M on primary fields}), 
(\ref{action of D on primary fields}), 
(\ref{action of K on primary fields})
において$x=0$とおくと、
\begin{align}
    &
    M_{μν}𝒪(0) = 𝒮_{μν}𝒪(0),
    \\ &
    D𝒪(0) = Δ𝒪(0),
    \\ &
    K_μ𝒪(0) = 0
    \label{another definition of primary fields}
\end{align}
が得られる。
特に
(\ref{another definition of primary fields})
はプライマリー場の定義とすることもできる。

この定義が元の定義を再現することを以下で示す。
$P_μ = ∂_μ$は常に成り立つので、
\footnote{
    並進によって内部空間の変換が引き起こされるような変な場は基本的に考えない。
}
$P_μ𝒪(x)=∂_μ𝒪(x)$であり、
\begin{align}
    𝒪(x) = \e^{x⋅P}𝒪(0)
\end{align}
と書ける。
Baker-Campbell-Hausdorffの公式
\begin{align}
    \e^{A}B\e^{-A} =
    \e^{[A,⋅]}B
    = B + [A,B] + \f{1}{2}[A,[A,B]]+⋯
\end{align}
を用いると、
\begin{align}
    M_{μν}𝒪(x)
    &
    = M_{μν}\e^{x⋅P}𝒪(0)
    \∅ &
    =\e^{x⋅P}(\e^{-x⋅P} M_{μν}\e^{x⋅P})𝒪(0)
    \∅ &
    =\e^{x⋅P} (M_{μν}-x_μP_ν+x_νP_μ)𝒪(0)
    \∅ &
    = (𝒮_{μν}-x_μ∂_ν+x_ν∂_μ)𝒪(x)
\end{align}
となるので、
(\ref{action of M on primary fields})
が得られる。
同様に
\begin{align}
    &
    \e^{-x⋅P}D\e^{x⋅P} = D + x⋅P
    \\ &
    \e^{-x⋅P}K_μ\e^{x⋅P}
    = K_μ + 2x_μD - 2x^νM_{μν}
    + 2 x_μ(x⋅P) - x²P_μ
\end{align}
を用い、
(\ref{another definition of primary fields})
を課すことで
(\ref{action of D on primary fields}), 
(\ref{action of K on primary fields})
が得られる。
微小変換に対する変換性から有限変換に対する変換性を構成できるので、(\ref{another definition of primary fields})
は
(\ref{definition of primary fields})
と等価である。

\subsection{
    ディセンダント場
}
プライマリー場$𝒪(x)$を微分して得られる場$P_μ𝒪(x),P_μP_ν𝒪(x),…$をディセンダント場と呼ぶ。
\begin{align}
    K_μP_ν𝒪(0)
    &
    = [K_μ,P_ν]𝒪(0)+P_νK_μ𝒪(0)
    \∅ &
    = (2Δδ_{μν}-2𝒮_{μν})𝒪(0)
\end{align}
が消えないことから、一般に$P_μ𝒪(x)$はプライマリー場ではない。
さらに、
\begin{align}
    DP_μ𝒪(x)
    &
    = (P_μD + [D,P_μ])𝒪(x)
    \∅ &
    = P_μ(D+1)𝒪(x)
    \∅ &
    = (x⋅∂ + Δ+1)(P_μ𝒪(x))
\end{align}
から、$P_μ𝒪(x)$は共形次元が1上がっていることが分かる。
$P_μ$を$n$回掛けたディセンダント場の共形次元は$Δ+n$となり、これをレベル$n$のディセンダント場という。
同様に、
\begin{align}
    DK_μ𝒪(x)
    &
    = (K_μD + [D,K_μ])𝒪(x)
    \∅ &
    = K_μ(D-1)𝒪(x)
    \∅ &
    = (x⋅∂ + Δ-1)(K_μ𝒪(x))
\end{align}
から$K_μ$は共形次元を1下げることが分かる。
しかし後で議論するように、ユニタリな場の理論では共形次元に下限が存在するので、$K_μ$を作用させていくと必ず$K_μ𝒪(0)=0$となる場が現れる。
これはプライマリー場に他ならない。
つまりプライマリー場は最低ウェイト状態であり、一般の場はプライマリー場に$P_μ$を作用させて構成できる。
プライマリー場$𝒪(x)$と、その微分で得られるディセンダント場をまとめて共形族(conformal family)と呼ぶ。
% また、似たような概念で、
% \begin{align}
%     𝒱 ≔ \Span\{P_μP_ν⋯𝒪(0)\}/\rm{permutation}
% \end{align}
% をVerma加群と呼ぶ。$/\rm{permutaion}$は$P$の置換についての同値類をとることを意味している。

\subsection{
    共形Casimir演算子
}
$d$次元共形代数は$d+2$次元Lorentz代数$𝔰𝔬(1,d+1)$に同型であることが知られている。
\footnote{
    $d$次元Minkowski時空における共形代数は$𝔰𝔬(2,d)$と同型である。
}
実際、
\begin{align}
    L_{μν} = M_{μν},
    \␣
    L_{-1μ} = \f{1}{2}(P_μ-K_μ),
    \␣
    L_{-10} = D,
    \␣
    L_{0μ} = \f{1}{2}(P_μ+K_μ)
\end{align}
によって$L_{ab}~(a,b=-1,0,…,d)$を定義し、$L_{ba} = -L_{ab}$を課すと、共形代数の交換関係は
\begin{align}
    [L_{ab},L_{cd}] = -η_{bc}L_{ad}+η_{ac}L_{bd}+η_{bd}L_{ac}-η_{bc}L_{ab}
\end{align}
という形にまとめられる。ここで$η$は符号$(-,+,…,+)$をもつMinkowski計量である。
これは$d+2$次元Lorentz代数$𝔰𝔬(1,d+1)$に同型である。
ここから共形Casimir演算子
\begin{align}\tcboxmath{
    L^2 = \f{1}{2}L_{ab}L^{ab}
    = \f{1}{2}M_{μν}M^{μν}-D²+\f{1}{2}(P_μK^μ+K^μP_μ)
}\end{align}
は共形代数の全ての演算子と交換する。
% \footnote{
%     これ以外にも以下のようなCasimir演算子がある。
%     \begin{align}
%         \f{1}{2}L_{ab}L^{bc}L_{cd}L^{da}
%     \end{align}
% }

ここで、
\begin{align}
    L^2 = \f{1}{2}M_{μν}M^{μν}-D(D-d)+P_μK^μ
\end{align}
と変形できる。
これをプライマリー場に作用させると、
\begin{align}
    L^2𝒪(0) = \Q(\f{1}{2}𝒮_{μν}𝒮^{μν}-Δ(Δ-d))𝒪(0)
\end{align}
となる。
$𝒪(0)$をスピン$l$のテンソル
\footnote{
    $l$階の対称トレースレステンソルのこと。この記事ではそのようなテンソルのみを扱う。
}
とすれば、$\f{1}{2}𝒮_{μν}𝒮^{μν}𝒪(0) = -l(l+d-2)𝒪(0)$となることが知られている
\footnote{
これを確かめる簡便な方法はLaplacianの極座標表示
\begin{align}
    ∂^2 = \f{1}{r^2}\Q[
        \Q(r\∂{r})^2 + dr\∂{r} + 𝓂^2
    ]
\end{align}
を用いる方法である。
\begin{align}
    ∂^2\f{1}{|y-x|^{d-2}}= 0
\end{align}
に注目し、$1/|y-x|^{d-2}$を$|x|/|y| ≕ r/|y|$のべきによって展開する。
$r^l$に比例する項は$C_l(𝒏)r^l$と書け、これに対し
\begin{align}
    -r^{l-2} 𝓂^2C_l(𝒏) = C_l(𝒏)∂^2r^l = l(l+d-2)C_l(𝒏)r^{l-2}
\end{align}
が成り立つ。ちなみに$C_l(𝒏)$はGegenbauer多項式と呼ばれる特殊関数になる。
$r^l$の係数はスピン$l$の表現に対応しており、この表現に対するCasimir演算子の固有値は$-l(l+d-2)$である。
}
ので、
\begin{align}
    L^2𝒪(0) = -λ_{Δ,l}𝒪(0),\␣ 
    λ_{Δ,l} = l(l+d-2) + Δ(Δ-d)
\end{align}
となる。
さらに、共形Casimir演算子が$P_μ$と交換することから、任意のディセンダント場はプライマリー場と同じ固有値を持つ。
% \subsection{
%     Weyl変換
% }
% 共形変換と似た概念として、Weyl変換がある。
% これは座標変換ではなく計量の変換であり、
% \begin{align}
%     g_{μν}(x) → Ω(x)^2g_{μν}(x)
% \end{align}
% というものである。
% 共形変換は平坦な空間の座標変換のうち、Weyl変換によって計量を元に戻せるものとして定義できる。
\end{document}